Arm 于 2011 年公布 Armv8 架构，并在随后发布 Cortex-A53、Cortex-A57 等首批
Armv8-A 处理器。Armv8-A 最初定义 AArch64 与 AArch32 两种执行状态，
但具体处理器可以只实现其中一部分；因此不能假设所有 Armv8-A 或后续 Armv9-A
处理器都能够运行 AArch32 软件。

本节以早期 Armv8-A 和 Cortex-A57 为教学基线，并参考
\emph{Arm Cortex-A Series Programmer's Guide for Armv8-A}\cite{armpg}。
后续架构版本已经加入 SVE/SVE2、MTE、Realm Management Extension 等能力，
实现相关软件时应同时查阅目标处理器对应版本的 Architecture Reference Manual 和勘误表。

\begin{description}
    \item[Large physical address] 更大的物理内存，使处理器能够访问超过 4GB 的物理内存。
    \item[64-bit virtual addressing] 使虚拟内存突破 4GB 限制。
    这对于使用内存映射文件 I/O 或稀疏寻址的现代桌面和服务器软件很重要。
    \item[Automatic event signaling] 可以实现节能、高性能的自旋锁。
    \item[Larger register files] 31 个 64 位通用寄存器可提高性能并减少堆栈使用。
    \item[Efficient 64-bit immediate generation] 降低对 literal pool 的需求。
    \item[PC-relative addressing] A64 提供多种不同范围的 PC 相对指令；例如
      \texttt{ADRP} 可在当前页附近的 $\pm 4$GB 范围生成页地址，
      \texttt{ADR} 和分支指令则具有各自更小的编码范围。
    \item[Additional 16KB and 64KB translation granules] 降低了翻译后备缓冲区 (TLB) 未命中率和页面遍历深度。
    \item[New exception model] 降低了操作系统和管理程序软件的复杂性。
    \item[Efficient cache management] 架构提供数据缓存清零与维护指令；
      EL0 是否能够执行部分缓存维护操作由更高异常级配置，不能默认开放。
    \item[Hardware-accelerated cryptography] 可选的密码扩展提供 AES、SHA 等指令，
      加速比取决于微架构、算法、数据规模和软件实现，不能使用固定倍数概括。
    \item[Load-Acquire, Store-Release instructions] 专为 C++11、C11、Java 内存模型而设计。
    它们能减少部分同步路径对显式屏障的需求，但不能替代所有屏障或语言级原子语义。
    \item[NEON double-precision floating-point advanced SIMD] 使得 SIMD 矢量化能够应用于更广泛的算法集，例如科学计算、高性能计算 (HPC) 和超级计算机。
\end{description}

\subsection{Cortex-A57 处理器}

现代 Arm Cortex 与 Neoverse 产品已经远超 Cortex-A57 所处的代际，
但 A57 是较早、资料完整的乱序 Armv8-A 核，适合用来理解经典 AArch64 实现。
本节讨论的是 Cortex-A57 的具体微架构，而不是所有 Armv8-A 处理器的通用属性。

Cortex-A57 面向移动和企业计算应用，可与 Cortex-A53 组成早期 Arm big.LITTLE 系统，
在性能与功耗之间进行任务迁移或异构调度。

Cortex‑A57 处理器具有与其他处理器的高速缓存一致性互操作性，包括用于 GPU 计算的 ARM Mali™ 系列图形处理单元 (GPU)，并为高性能企业应用程序提供可选的可靠性和可扩展性功能。
它提供了比 ARMv7 架构的 Cortex‑A15 处理器更高的性能，并具有更高的能效水平。
可选密码扩展可显著加速部分算法，但实际收益依赖具体实现与工作负载。

\begin{figure}[htbp]
  \centering
  \resizebox{0.72\textwidth}{!}{%
    \begin{tikzpicture}[
      font=\sffamily\small,
      block/.style={draw=booknavy, rounded corners=2mm, fill=bookpaper,
        align=center, minimum height=0.72cm},
      unit/.style={block, fill=white, minimum height=1.15cm},
      cluster/.style={draw=booknavy, rounded corners=2mm, fill=bookline,
        minimum width=8.4cm, minimum height=5.0cm}
    ]
      \fill[bookblue!55, rounded corners=3mm] (0,0) rectangle (10.2,10.9);
      \node[block, minimum width=9.5cm] at (5.1,10.35)
        {CoreSight multicore debug and trace};
      \node[block, minimum width=9.5cm] at (5.1,9.45)
        {Generic Interrupt Controller};

      \node[cluster] at (5.75,6.55) {};
      \node[cluster] at (5.4,6.25) {};
      \node[cluster] at (5.05,5.95) {};
      \node[cluster, fill=bookpaper!80!bookline] at (4.7,5.65) {};
      \node[font=\sffamily\scriptsize\bfseries, anchor=east] at (8.72,4.05)
        {Cores 0--3};

      \node[unit, minimum width=4.6cm, minimum height=2.75cm] at (3.25,6.75)
        {Cortex-A57\\out-of-order core};
      \node[unit, minimum width=2.25cm, minimum height=1.3cm] at (6.9,7.5)
        {Advanced SIMD\\and crypto};
      \node[unit, minimum width=2.25cm, minimum height=1.05cm] at (6.9,6.0)
        {Floating-point\\unit};
      \node[unit, minimum width=2.25cm, minimum height=1.05cm] at (6.9,4.65)
        {Memory\\protection unit};
      \node[unit, minimum width=2.15cm, minimum height=1.05cm] at (1.95,4.65)
        {L1 instruction\\cache};
      \node[unit, minimum width=2.15cm, minimum height=1.05cm] at (4.3,4.65)
        {L1 data cache\\with ECC};
      \node[unit, minimum width=4.5cm, minimum height=0.75cm] at (3.25,3.55)
        {Performance Monitor Unit};

      \node[block, minimum width=1.9cm] at (1.45,2.25) {SCU};
      \node[block, minimum width=1.9cm] at (3.65,2.25) {ACP};
      \node[block, minimum width=5.0cm] at (7.3,2.25)
        {Integrated shared L2 cache with ECC};
      \node[block, minimum width=9.5cm, minimum height=0.8cm] at (5.1,0.85)
        {AMBA 4 ACE coherent interface};
    \end{tikzpicture}%
  }
  \caption{Cortex-A57 集群结构（教学简图）}
  \label{fig:cortex-a57}
\end{figure}

Cortex-A57 实现 Armv8.0-A，并支持 AArch64 和 AArch32 执行状态。
它支持多核操作，在单个集群中具有一到四核多处理器。
该核采用 AMBA 4 ACE 一致性接口；系统级互连是否使用其他协议取决于具体 SoC。
CoreSight 组件可用于调试和跟踪。

Cortex‑A57 处理器具有以下特性：

\begin{itemize}
  \item 乱序，多于 15 级流水线。
  \item 省电功能包括路径预测、标记缩减和缓存查找抑制。
  \item 通过重复执行资源增加峰值指令吞吐量。
    具有本地化解码、3 个宽解码带宽的功率优化指令解码。
  \item 性能优化的 L2 缓存设计使集群中的多个核心可以同时访问 L2。
\end{itemize}

\input{armv8/armv8-basic}
\input{armv8/regs}
\input{armv8/isa-in-nutshell}
\input{armv8/a64-isa}
\input{armv8/float-neon}
\input{armv8/abi}
\input{armv8/exception}
\input{armv8/cache}
\input{armv8/mmu}
\input{armv8/memory-ordering}
\input{armv8/multi-core-processors}
\input{armv8/power-management}
\input{armv8/big-little-tec}
\input{armv8/security}
\input{armv8/debug}
\subsection{Armv8 虚拟平台与模型}
\label{sec:armv8-virtual-platforms}

虚拟平台允许固件、Bootloader、内核和部分驱动在目标芯片可用前开始开发，也适合把启动回归和异常路径纳入持续集成。它解决的是“软件在给定体系结构和平台契约下是否正确”，并不天然回答真实芯片的周期级性能、电气时序、模拟行为或制造差异。

Arm Fast Models 是构建高速度功能模型的组件与工具体系，Fixed Virtual Platform（FVP）则是由处理器、互连、内存和外设模型组成的可直接运行平台\cite{arm-fast-models,arm-fvp-overview}。不同版本的产品名称、可配置组件和命令行参数会变化，项目应固定所用模型版本并以随版本发布的参考手册为准。

\subsubsection{先区分模型精度}

“能够启动 Linux”只证明一部分功能契约成立。选择模型前应区分以下精度维度：

\begin{description}
  \item[指令精度] 指令结果、异常和系统寄存器行为是否符合相应架构版本；
  \item[平台精度] 内存映射、中断控制器、定时器、UART、PCIe 等设备是否与目标平台契约一致；
  \item[事务精度] 总线请求、顺序和争用是否以足够细节建模；
  \item[时间精度] 指令、设备和互连延迟能否代表真实硬件；
  \item[电气精度] 电压、边沿、抖动、模拟前端和信号完整性是否被建模。
\end{description}

常用虚拟平台通常重视指令和软件可见的平台行为，对周期、总线争用和电气行为作出简化。模型报告的虚拟时间适合驱动定时器和超时状态机，却不能直接作为产品 WCET、启动时间或吞吐量结论。

\begin{table}[htbp]
  \centering
  \caption{嵌入式开发环境的典型适用范围}
  \begin{tabular}{p{0.18\textwidth}p{0.34\textwidth}p{0.36\textwidth}}
    \toprule
    环境 & 适合验证 & 不能单独证明 \\
    \midrule
    主机单元测试 & 算法、解析器、状态机、错误处理 & 目标 ISA、MMIO、中断和真实并发 \\
    ISA 模拟器 & 指令语义、ABI、异常入口 & 完整 SoC 设备契约与时序 \\
    QEMU 系统模拟 & 操作系统启动、通用虚拟设备、自动化回归 & 目标专用外设、电气行为和精确性能 \\
    \shortstack[l]{FVP/Fast\\Models} & Arm 架构功能、平台固件、GIC、PSCI、调试流程 & 真实芯片周期、PHY、模拟和板级问题 \\
    RTL/仿真加速器 & 选定硬件实现和寄存器级交互 & 完整产品环境及未建模器件 \\
    FPGA 原型 & 接近硬件的接口与长时间软件运行 & 最终频率、功耗、模拟和硅片勘误 \\
    目标板/HIL & 电气、性能、功耗、温度与完整系统行为 & 未覆盖样本之外的绝对正确性 \\
    \bottomrule
  \end{tabular}
\end{table}

项目不应寻找一个“最真实”的环境替代所有其他环境，而应让同一组可移植测试在多个层次运行，并明确每层新增的证据。

\subsubsection{Fast Models 与 FVP}

Fast Models 通过高层组件和事务级连接描述处理器与平台。模型通常采用指令翻译和时间量化提高速度，因此可在普通开发主机上运行完整固件或操作系统。FVP 将一组模型、拓扑和默认参数封装为可执行平台，常同时提供串口、网络、块设备、显示和调试连接。

虚拟平台适合：

\begin{itemize}
  \item 验证 EL3、EL2、EL1 与 EL0 的启动和异常级交接；
  \item 开发 Trusted Firmware-A、UEFI、U-Boot、内核和初始根文件系统；
  \item 验证 GIC、通用定时器、PSCI、SMCCC、SMMU 等架构接口；
  \item 执行确定性的启动回归、故障注入和多配置测试；
  \item 在目标板数量有限时为开发者和 CI 提供并行环境。
\end{itemize}

虚拟平台不能替代：

\begin{itemize}
  \item DDR 训练、PHY 校准、SerDes、射频和模拟外设验证；
  \item cache、互连、DRAM 与 DMA 争用下的真实性能测量；
  \item 依赖精确中断到达、总线 backpressure 或器件内部延迟的实时结论；
  \item 电源时序、复位毛刺、热行为、EMC 和板级信号完整性测试；
  \item 目标芯片未公开实现细节、勘误或生产批次差异的复现。
\end{itemize}

模型与硬件结果不一致时，应先判断差异属于架构允许的实现选择、模型缺失、软件依赖未定义行为，还是目标芯片勘误。不能默认把任一方当作绝对真值。

\subsubsection{Architecture/基础平台}

面向架构的软件通常应先在通用 Architecture Model 或历史上称为 Foundation Platform 的环境运行。这类平台使用架构模型处理器，不对应某一颗具体 Cortex 内核，平台设备也以支持软件开发所需的最小集合为主。其价值是减少具体 SoC 特性的干扰，使问题集中在架构契约。

适合在基础平台验证的内容包括：

\begin{itemize}
  \item AArch64 复位入口、异常向量和页表建立；
  \item EL3 到非安全世界的交接、PSCI CPU 启停和系统复位；
  \item GIC 初始化、IPI、通用定时器和多核 bring-up；
  \item 内核启动协议、Device Tree、initramfs 与早期控制台；
  \item 对架构可选特性的运行时探测与拒绝路径。
\end{itemize}

软件不得根据模型名称猜测功能。应通过 ID 寄存器、固件表或 Device Tree 探测处理器扩展、地址宽度、GIC 版本和设备实例。测试某项可选扩展时，应显式记录模型参数与探测结果；否则模型升级后默认值变化可能让回归失去可比性。

基础平台的设备树只能描述该虚拟平台，不能直接复制到产品 SoC。尤其要核对：

\begin{itemize}
  \item RAM、固件保留区和 MMIO 地址；
  \item GIC distributor、redistributor 和 ITS 布局；
  \item UART 类型、时钟和中断号；
  \item PSCI conduit 是 SMC 还是 HVC；
  \item PCIe ECAM、DMA coherency 与 IOMMU 拓扑；
  \item CPU 数量、MPIDR 拓扑和 NUMA/cluster 描述。
\end{itemize}

基础平台通过不代表产品 BSP 已完成，但它可以证明大部分架构相关代码未被某个开发板的偶然行为掩盖。

\subsubsection{Base Platform FVP}

Base Platform FVP 提供比最小基础平台更完整、可配置的系统拓扑，常用于平台固件、操作系统和参考软件栈的联合开发。具体发行版可能允许选择架构模型或特定处理器模型，并配置核数、内存、GIC、SMMU、PCIe、VirtIO 和显示等组件；可用参数必须通过该版本的帮助和参数列表确认。

Base FVP 更适合验证：

\begin{itemize}
  \item 多阶段安全启动以及 BL1、BL2、BL31、UEFI/U-Boot 的镜像交接；
  \item 多 cluster 拓扑、CPU hotplug、低功耗接口和二级启动；
  \item EL2 hypervisor、Stage-2 地址转换和虚拟中断；
  \item SMMU、PCIe 与 DMA 隔离的软件配置；
  \item OP-TEE 等安全世界服务与普通世界客户端的协议；
  \item 大型 Linux 发行镜像和参考平台集成测试。
\end{itemize}

如果目标产品使用厂商专用 FVP，应把通用 Base FVP 与厂商模型分成两条测试线：前者保护架构和上游兼容性，后者覆盖专用地址图、控制器与固件。只在厂商模型运行会让软件不必要地依赖私有行为；只在通用模型运行又会遗漏真实 SoC 契约。

\subsubsection{启动镜像与平台契约}

一次可复现的虚拟平台启动至少要冻结以下输入：

\begin{itemize}
  \item FVP 可执行文件、组件包和许可证版本；
  \item 平台参数及其非默认值；
  \item 所有固件、Bootloader、内核、DTB、initramfs 和磁盘镜像摘要；
  \item 每个镜像的加载地址、入口地址和安全属性；
  \item 串口、网络、块设备和半主机接口配置；
  \item 测试超时采用虚拟时间、主机时间还是二者组合。
\end{itemize}

建议把命令行参数固化为受版本控制的配置文件，并在运行前导出平台参数清单。以下命令只表达调用模式，选项和参数名必须按所用 FVP 版本替换：

\begin{lstlisting}[language=bash,caption={固定虚拟平台运行输入的示意流程}]
${FVP_BIN} --version
${FVP_BIN} --list-params > model-params.txt
${FVP_BIN} -f ci/fvp/base.cfg \
  --data ${BOOT_IMAGE}@${BOOT_ADDRESS}
\end{lstlisting}

CI 不应依赖交互式终端窗口。串口输出应写入文件或标准流，网络端口应动态分配或隔离，测试完成则由半主机退出、模型控制接口或稳定的串口协议返回机器可读结果。单纯搜索“login:”只能证明启动到某一阶段，不能证明所有服务和设备已经满足健康条件。

\subsubsection{Device Tree、ACPI 与固件接口}

虚拟平台与客户机之间的核心不是启动命令，而是平台描述和固件 ABI。Device Tree/ACPI、PSCI、SMCCC、UEFI memory map 和中断拓扑必须相互一致。

常见错误包括：

\begin{itemize}
  \item DTB 来自另一版本平台，导致 UART 地址或 GIC 区域不匹配；
  \item 固件保留内存未从操作系统内存图排除；
  \item 设备声明 DMA coherent，但模型或软件路径并未实现相应语义；
  \item PSCI conduit、function ID 或异常级路由不一致；
  \item 内核启用较新架构特性，而固件保存/恢复上下文不完整；
  \item 虚拟设备被当作产品硬件，驱动接口在移植时被错误固化。
\end{itemize}

每次平台升级都应先比较参数、内存图、DT/ACPI 和启动日志，再解释测试行为变化。只更新模型二进制而沿用旧基线，可能把平台变化误判为软件回归。

\subsubsection{调试与可观测性}

Fast Models/FVP 通常提供调试接口、组件状态、事件日志和统计信息；可接入 Arm 调试器，部分平台也能暴露 GDB 兼容连接。调试能力随模型和许可变化，项目应把最小无 GUI 调试路径写入仓库。

推荐保留四类证据：

\begin{enumerate}
  \item 从复位开始的全部串口输出；
  \item 模型版本、参数和镜像摘要；
  \item 失败时的 PC、异常级、关键系统寄存器和调用栈；
  \item 能够关联固件、内核与用户空间阶段的稳定事件 ID。
\end{enumerate}

半主机文件访问和模型专用退出设备很适合测试自动化，但必须与量产路径隔离。若代码在无意中依赖半主机，目标板上可能直接异常或永久等待调试器。

断点会停止一个或多个虚拟处理器，但计时器、外设和其他核是否同时暂停由模型配置决定。分析多核竞态时，应优先使用事件 trace 和确定性重放，避免把交互式单步产生的新时序当成原故障。

\subsubsection{自动化、确定性与差分验证}

同一模型的确定性执行有助于复现，但主机线程调度、网络桥接、随机种子和异步 I/O 仍可能引入变化。测试应固定随机种子和输入，禁用不需要的主机集成，并同时设置虚拟时间上限与墙上时间上限，防止模型无进展时占满 CI worker。

建议建立以下差分测试：

\begin{itemize}
  \item 同一镜像在 FVP、QEMU 与目标板运行架构自检和启动 smoke test；
  \item 同一协议状态机在主机、固件单元测试和完整系统中使用同一测试向量；
  \item 对系统寄存器、异常综合征和页表结果比较语义，不比较无保证的周期值；
  \item 对可选特性分别测试“实现”和“未实现”配置，验证运行时探测；
  \item 将模型专用设备从产品配置中移除后再次构建，防止依赖泄漏。
\end{itemize}

差分结果不应要求每一条日志字节完全相同。并行启动顺序、设备枚举号和地址随机化可能合法变化；测试 oracle 应围绕 ABI、不变量、错误码、状态转移和资源上限设计。

\subsubsection{从模型到目标硬件的退出门}

虚拟平台通过后，至少还要在 RTL/FPGA 原型或目标板完成以下验证：

\begin{itemize}
  \item 真实时钟、复位、电源域和 DDR 初始化；
  \item cache 一致性、DMA、IOMMU 和共享互连压力；
  \item IRQ 延迟、实时抖动、启动时间和吞吐量；
  \item PHY、存储介质、网络、电源管理和热状态切换；
  \item 芯片勘误规避、边界温度电压以及长稳运行；
  \item 调试关闭、Secure Boot 和量产生命周期状态。
\end{itemize}

模型测试失败通常定位软件契约，硬件测试失败还可能来自电气和实现问题。两类证据结合，才能避免在昂贵目标板上重复发现纯软件错误，也避免用虚拟平台的成功掩盖真实物理风险。

\subsubsection{虚拟平台检查清单}

\begin{itemize}
  \item 是否说明所需的是指令、平台、事务、时间还是电气精度？
  \item 模型、组件、参数、镜像和平台描述是否全部可追溯？
  \item 软件是否通过架构机制探测可选特性，而不是根据模型名称猜测？
  \item Device Tree/ACPI、内存图、GIC、PSCI 和异常级契约是否一致？
  \item CI 是否无 GUI、具有双重超时并能区分测试失败与模型基础设施失败？
  \item 测试 oracle 是否验证稳定语义，而不是无保证的日志顺序和周期数？
  \item 是否明确列出必须在 FPGA/目标板/HIL 上重新验证的结论？
\end{itemize}

